\chapter{运动想象脑电解码与端侧系统关键技术基础}
\label{cha:background}

\section{运动想象脑电信号原理}
\subsection{运动想象脑电信号特性}

运动想象（Motor Imagery, MI）是指个体在不产生实际肢体运动输出的前提下，在脑内主动表征某一动作过程的心理活动。尽管外周肌肉并未完成真正收缩，但感觉运动皮层仍会出现与运动准备和运动执行相近的节律调制现象，因此该范式成为非侵入式脑机接口中最常用的控制任务之一。对经典左右手运动想象而言，判别信息主要分布在中央区手部投影附近，其中 C3 与 C4 通常承担最显著的侧化信息，而 Cz 则在中央 $\beta$ 节律或部分个体的判别模式中提供补充作用 \cite{Pfu06,Mcf04}。因此，从生理基础看，运动想象脑电并不是任意通道上的弱相关波动，而是围绕感觉运动皮层局部区域展开的、具有明确空间指向性的节律响应。

这种空间局部性在多项经典研究中得到了直接支持。Pfurtscheller 等在 $9$ 名健康右利手受试者、$60$ 导联条件下观察到，手部运动想象会在 C3 和 C4 附近诱发以 $10\ \text{Hz}$ 和 $20\ \text{Hz}$ 成分为代表的广带事件相关去同步化（Event-Related Desynchronization, ERD），并指出最具判别价值的电极位置主要集中在 C3、Cz 和 C4 \cite{Pfu06}。McFarland 等在 $33$ 名受试者、$64$ 导联实验中的拓扑分析进一步表明，手部运动想象期间中央区节律变化具有稳定可观测性，但不同频带的峰值位置并不完全重合 \cite{Mcf04}。这些证据共同说明，左右手运动想象的有效解码信息主要来源于中央手区附近的局部时空模式，而不能简单压缩为一对固定电极之间的镜像差异。

从频谱特征看，运动想象脑电最关键的两类节律是 $\mu$ 节律与 $\beta$ 节律。前者通常位于 $8$--$13\ \text{Hz}$ 左右，后者常位于 $13$--$30\ \text{Hz}$ 范围内；两者都与感觉运动皮层状态调节相关，但在空间分布和功能表达上并不等价。Pfurtscheller 等报告最具反应性的 $\mu$ 成分主要分布于 $9$--$14\ \text{Hz}$，其平均反应频率为 $11.7 \pm 0.4\ \text{Hz}$，而去同步化中心频率约为 $10.9 \pm 0.9\ \text{Hz}$ \cite{Pfu06}。McFarland 等则将 $\mu$ 频段定义为 $8$--$12\ \text{Hz}$、$\beta$ 频段定义为 $18$--$25\ \text{Hz}$，并发现二者在运动想象中的拓扑中心并不一致：$\mu$ 节律更偏向双侧感觉运动区，而 $\beta$ 节律则更倾向于中央或额中央邻近区域 \cite{Mcf04}。因此，在运动想象脑电分析中，将 $\mu$ 与 $\beta$ 简化为同一种“SMR 宽带特征”会掩盖其潜在的功能差异，也不利于后续模型对不同频带信息的针对性利用。

除了群体层面的共性之外，运动想象脑电还表现出明显的个体差异。Zapała 等对 $20$ 名右利手与 $20$ 名左利手受试者的比较表明，右利手在左手运动想象任务中表现出更强的 $\alpha/\mu$ 去同步化，右顶叶感觉运动相关簇中的平均功率变化分别为 $-2.59\ \text{dB}$ 和 $-1.46\ \text{dB}$，差异具有统计学意义（$p=0.031$）；对应的单次在线控制准确率分别为 $51.54\%\pm16.43$ 与 $43.25\%\pm8.24$ \cite{Zap20b}。这说明，左右手运动想象虽然具有稳定的中央区节律基础，但其侧化强度并非在所有个体中都呈现同样的镜像结构。对论文后续的解码任务而言，这一事实意味着通道权重、时间窗位置和频带响应都应允许保留一定的个体适配空间，而不宜预设完全刚性的对称先验。

\subsection{脑电信号预处理技术}

原始脑电信号幅值通常处于微伏量级，且易受工频干扰、眼动伪迹、肌电噪声、接触阻抗变化以及基线漂移等因素影响。因此，在进入特征提取与模型建模阶段之前，必须通过预处理步骤提升有效神经信息的信噪比。对运动想象任务而言，常见的预处理流程包括带通滤波与陷波滤波、坏段剔除、重参考、按提示事件切片，以及在必要时结合独立成分分析或空间滤波方法抑制伪迹和无关活动 \cite{lotte2018review, ang2012filter}。这些步骤的目的并不在于人为制造判别性，而是在尽可能保留感觉运动节律特征的前提下，削弱与任务无关的非平稳干扰。

然而，运动想象脑电中的预处理并非完全中性的“前置清洗”过程。时间窗和基线的选择会直接影响 ERD/ERS 的可见形态及其后续可分性。以时间窗为例，Xu 等在单试次 $\mu$ 节律抑制检测研究中比较了不同窗口长度和不同潜伏期设置，发现采用群体平均潜伏期时，运动想象检测成功率约为 $58.47\%$；若引入个体化潜伏期搜索，成功率可提升至 $74.14\%$，差异显著（$t(11)=4.766,\, p=0.001$）\cite{Xu19}。这说明，$635\ \text{ms}$ 的群体平均潜伏期和 $150\ \text{ms}$ 的短窗口并不能被简单理解为对所有受试者都普适有效的固定参数，而应被视为“短时局部检测需要嵌入更宽搜索范围”的经验依据。换言之，运动想象信号的时间表达具有明显的个体依赖性，若在预处理阶段直接套用单一窗口设定，往往会损失部分真实有效的节律信息。

基线窗口的选择同样会改变对运动想象脑电的解释。Rimbert 等在 $71$ 名健康右利手、运动想象初学者数据上的分析表明，使用不同的基线区间会显著改变 ERD 的空间分布和强度估计；相较于更靠近下一试次提示或更易混入前一试次反弹效应的基线，较早的基线设置在左右手 MI-vs-rest 任务中可带来约 $4\%$ 至 $8\%$ 的分类表现差异 \cite{Rim23b}。这一结果说明，基线并不是一个与生理解释无关的技术细节。如果基线区间被前一试次残留的同步化反弹或临近提示的预期活动污染，那么后续观察到的 ERD 幅度就可能被高估或低估，最终影响对模型性能和信号机制的判断。

综合上述研究可以发现，运动想象脑电的预处理策略应当服从其生理特征，而不是脱离任务机制进行统一模板化设置。具体而言，中央区局部通道应被优先保留，$\mu$ 与 $\beta$ 节律不宜被粗略混合为单一宽频带特征，时间窗选择需要考虑个体潜伏期差异，基线区间则应尽量避开前后试次的污染区间 \cite{Pfu06,Mcf04,Xu19,Rim23b}。这也为后续解码模型的设计提供了明确约束：只有在预处理阶段尽可能忠实保留感觉运动节律的空间、频谱与时间结构，后续的特征提取与注意力建模才能建立在可靠的生理基础之上。

\section{自注意力机制与Transformer架构}

\subsection{Transformer 模型基本原理}
Transformer 的核心思想是在不依赖递归结构的前提下直接建模序列内部的成对关系 \cite{vaswani2017attention}。从结构组成来看，该模型通常由输入嵌入、位置编码、若干编码器块以及任务相关输出头构成。对于长度为 $N$ 的输入序列，设其嵌入表示为 $\mathbf{E} \in \mathbb{R}^{N \times d}$，位置编码矩阵为 $\mathbf{P} \in \mathbb{R}^{N \times d}$，则进入编码器之前的初始表示可写为：
\begin{equation}
\mathbf{Z}^{(0)} = \mathbf{E} + \mathbf{P},
\label{eq:input_embed}
\end{equation}
其中 $d$ 表示嵌入维度。通过上述表示方式，模型在保留原始特征内容的同时引入了序列位置相关信息，从而为后续注意力计算提供顺序依据。

位置编码是 Transformer 区别于无序集合建模的重要组成部分。在经典 Transformer 中，位置编码常采用正弦余弦形式，其定义如下：
\begin{equation}
\begin{aligned}
\mathbf{P}_{(pos,2i)} &= \sin \left(\frac{pos}{10000^{2i/d}}\right), \\
\mathbf{P}_{(pos,2i+1)} &= \cos \left(\frac{pos}{10000^{2i/d}}\right),
\end{aligned}
\label{eq:position_encoding}
\end{equation}
式中，$pos$ 表示序列位置索引，$i$ 表示特征维度索引。式 \eqref{eq:position_encoding} 使模型能够在不引入额外递归结构的前提下感知顺序关系，这对脑电时序信号的上下文建模具有基础性作用。

编码器块决定了 Transformer 的核心特征提取过程。其内部通常包含多头自注意力（Multi-Head Attention, MHA）子层和前馈网络（Feed-Forward Network, FFN）子层，并在每个子层外部引入残差连接与层归一化（Layer Normalization, LN）。对于第 $l$ 层编码器，其输出可表示为：
\begin{equation}
\begin{aligned}
\hat{\mathbf{Z}}^{(l)} &= \text{MHA}(\text{LN}(\mathbf{Z}^{(l-1)})) + \mathbf{Z}^{(l-1)}, \\
\mathbf{Z}^{(l)} &= \text{FFN}(\text{LN}(\hat{\mathbf{Z}}^{(l)})) + \hat{\mathbf{Z}}^{(l)}.
\end{aligned}
\label{eq:encoder_block}
\end{equation}
其中，$\hat{\mathbf{Z}}^{(l)}$ 表示注意力子层输出，$\mathbf{Z}^{(l)}$ 表示第 $l$ 层编码器的最终输出。残差连接有助于稳定深层网络训练，层归一化则用于缓解不同样本间特征分布波动，因此二者已成为 Transformer 编码器中的标准配置。

在脑电分类任务中，Transformer 往往只保留编码器结构，而不使用原始序列到序列任务中的解码器分支。这是因为运动想象、情绪识别和癫痫检测等任务的目标通常是对输入脑电片段进行类别判定，而非生成新的输出序列。基于这一任务特征，模型一般在若干编码器层之后接入池化层或全连接分类头，将全局时序表征映射到具体类别空间 \cite{songEEG2022, xie_transformer-based_2022}。

\subsection{高效自注意力机制}

在上述基本架构中，自注意力层既决定了 Transformer 的表征能力，也决定了其主要计算代价。对高频采样、长时间窗输入的脑电数据而言，这一部分的复杂度分析尤其关键，因为它直接关系到模型能否在边缘硬件上完成实时推理。

\subsubsection{标准自注意力机制与计算复杂度瓶颈}

给定经过预处理和通道映射的输入连续脑电特征序列 $\mathbf{X} \in \mathbb{R}^{N \times d}$，其中 $N$ 为序列时间步长度，$d$ 为数据的特征维度。标准的自注意力机制首先通过三个不同的可学习参数矩阵 $\mathbf{W}_Q, \mathbf{W}_K, \mathbf{W}_V \in \mathbb{R}^{d \times d_k}$ 对输入 $\mathbf{X}$ 进行线性映射，分别得到查询矩阵（Query, $\mathbf{Q}$）、键矩阵（Key, $\mathbf{K}$）以及值矩阵（Value, $\mathbf{V}$），其计算过程如式 \eqref{eq:qkv} 所示：
\begin{equation} \label{eq:qkv}
\mathbf{Q} = \mathbf{X}\mathbf{W}_Q, \quad \mathbf{K} = \mathbf{X}\mathbf{W}_K, \quad \mathbf{V} = \mathbf{X}\mathbf{W}_V
\end{equation}

随后，算法通过计算 $\mathbf{Q}$ 和 $\mathbf{K}$ 之间的点积来衡量序列中不同位置的相关性，并在除以缩放因子 $\sqrt{d_k}$ 后经过 Softmax 激活函数得到注意力权重分数矩阵。最终的结果输出即是对 $\mathbf{V}$ 的加权求和，定义如式 \eqref{eq:attention} 所示：
\begin{equation} \label{eq:attention}
\text{Attention}(\mathbf{Q}, \mathbf{K}, \mathbf{V}) = \text{Softmax}\left(\frac{\mathbf{Q}\mathbf{K}^T}{\sqrt{d_k}}\right) \mathbf{V}
\end{equation}

标准自注意力的主要瓶颈来自相关性矩阵的构造过程。计算 $\mathbf{Q}\mathbf{K}^T$ 需要枚举序列中任意两个位置之间的相互作用，因此会生成一个尺寸为 $N \times N$ 的稠密注意力矩阵。由此带来的存储复杂度为 $O(N^2)$，主要计算项约为 $O(N^2 \cdot d)$。当脑电分析采用较长时间窗或较高采样率时，序列长度 $N$ 会迅速增大，进而导致注意力矩阵的存储与乘法运算成本同步上升。这也是标准全局自注意力在资源受限平台上部署时面临的主要瓶颈。

\subsubsection{局部非重叠与滑动窗口自注意力架构}

为了降低全局注意力带来的二次复杂度，局部窗口注意力将计算范围收缩到固定长度的局部邻域内，而不再令整个序列中的所有位置两两交互 \cite{liu2021swin}。

其复杂度优势可由窗口划分方式加以说明。设长度为 $N$ 的输入序列被划分为 $\frac{N}{M}$ 个互不重叠的局部窗口，其中每个窗口包含 $M$ 个特征片段，则注意力计算仅在单个窗口内部执行。此时，每个窗口的计算代价约为 $O(M^2 \cdot d)$，整体复杂度可写为：
\begin{equation} \label{eq:complexity_local}
\Omega(\text{W-MSA}) = \frac{N}{M} \cdot O(M^2 \cdot d) = O(N \cdot M \cdot d)
\end{equation}
当窗口长度 $M$ 被视为预设常数且满足 $M \ll N$ 时，式 \eqref{eq:complexity_local} 表明该类方法的复杂度可近似看作关于序列长度 $N$ 的线性增长。因此，局部窗口注意力在保留时序局部建模能力的同时，能够显著缓解长序列脑电输入带来的存储与计算压力。

但仅采用非重叠窗口仍难充分保持时序连续性，因为窗口边界会截断跨片段的信息交互。基于这一问题，许多改进模型进一步引入滑动窗口或移位窗口机制 \cite{liu2021swin, hua2023improved}。其基本做法是在相邻两层之间对窗口划分边界进行固定步长平移，使原本属于不同窗口的特征片段在下一层中进入同一局部感受域，从而建立跨窗口信息交换通路。借助这一设计，模型能够在保持较低复杂度的同时逐步扩大有效感受野，实现由局部依赖到更大范围上下文关系的层级建模。

脑电任务中的相关研究进一步验证了上述结论。Wan 等将通道注意力与 Swin 风格窗口机制结合，用于运动想象脑电的时空特征提取 \cite{wan2023novel}；Hua 等提出的 CSANet 则在卷积特征提取之后引入滑动窗口与两级注意力结构，以增强多尺度时空信息建模能力 \cite{hua2023improved}。这些代表性工作表明，局部窗口注意力并不仅是降低复杂度的技术手段，同时也是适配脑电局部节律模式与非平稳特征的有效结构设计。

\section{端侧部署、传输与采集关键技术}

\subsection{边缘部署的核心约束指标}

面向便携式脑机接口的端侧部署，不仅要关注分类精度，还必须同时满足推理时延、片上存储、瞬时功耗与单次推理能耗等工程约束。若模型总参数量为 $N_p$，权重量化位宽为 $b$ bit，则仅权重存储开销可近似表示为
\begin{equation}
S_w = \frac{N_p \cdot b}{8},
\label{eq:edge_weight_storage}
\end{equation}
式中，$S_w$ 的单位为字节（Byte）。除权重之外，特征图缓存、算子中间张量以及运行时缓冲区同样会占用片上 SRAM 或外部存储空间，因此实际部署中的内存压力通常高于单纯由参数量估计得到的结果。另一方面，若单次推理期间的平均功率为 $P_{\mathrm{avg}}$，推理时延为 $T_{\mathrm{inf}}$，则单次推理能耗可写为
\begin{equation}
E_{\mathrm{inf}} = P_{\mathrm{avg}} \cdot T_{\mathrm{inf}}.
\label{eq:edge_energy}
\end{equation}
因此，对实时脑机接口而言，模型是否适合端侧运行并不由某一项指标单独决定，而取决于精度、时延、存储与能耗之间的综合平衡。

已有运动想象脑电部署研究说明，端侧可行性的第一性问题并非“能否运行”，而是“在多小的资源预算下仍能保持可接受性能”。Wang 等将 EEGNet 变体部署到 Cortex-M4F 与 Cortex-M7 微控制器上，报告标准模型在二分类、三分类和四分类任务上分别达到 $82.43\%$、$75.07\%$ 和 $65.07\%$ 的准确率；缩放后的轻量版本在内存下降低 $7.6$ 倍时仅损失约 $0.31\%$ 准确率，最小模型在 Cortex-M4F 上的单次推理时延约为 $101\ \text{ms}$、能耗约为 $4.28\ \text{mJ}$ \cite{Wan20}。Schneider 等在 PULP 架构的 Mr.\ Wolf 平台上进一步给出量化部署结果：其 8-bit 量化 EEGNet 在仅带来约 $0.4\%$ 精度损失的情况下，实现了 $5.82\ \text{ms}$ 的单次推理时延和 $0.627\ \text{mJ}$ 的单次能耗 \cite{Sch20}。在更低成本的 Arduino Nano 33 Sense 上，En\'eriz 等同样验证了训练后量化（PTQ）对 MI 脑电模型的适用性，平均精度下降约为 $2.64 \pm 0.77\%$，推理时延约为 $137\ \text{ms}$，单次能耗约为 $2.55\ \text{mJ}$ \cite{Ene23}。这些结果共同说明，紧凑型 MI 脑电解码网络已经能够在比 RK3568 更弱的硬件平台上完成实时或准实时推理，因此边缘部署在原理上是成立的。

\subsection{模型压缩、量化与算子调度}

在边缘设备上部署深度模型时，最常见的优化路径包括结构剪枝、低比特量化以及面向特定算子的调度优化。剪枝的目标在于移除冗余通道、卷积核或连接权重，从而降低参数量与乘加运算次数；量化则将浮点权重和激活映射到低位宽整数表示，以减少存储占用并提升整数运算单元的利用效率。对对称量化而言，若浮点张量中的元素为 $x$，缩放因子为 $s$，则量化与反量化过程可写为
\begin{equation}
\begin{aligned}
q &= \mathrm{clip}\!\left(\mathrm{round}\!\left(\frac{x}{s}\right), q_{\min}, q_{\max}\right), \\
\hat{x} &= s \cdot q,
\end{aligned}
\label{eq:symmetric_quant}
\end{equation}
其中 $q$ 为整数域表示，$\hat{x}$ 为反量化后的近似值，$q_{\min}$ 与 $q_{\max}$ 由目标位宽决定。当 $b=8$ 时，式 \eqref{eq:symmetric_quant} 可直接对应 INT8 推理流程。量化的本质是在数值精度和硬件效率之间做受控折中：位宽越低，存储与带宽开销越小，但过度压缩会引入累积误差，导致判别边界退化。

现有文献对这一折中的可控性给出了较明确的工程证据。Schneider 等的量化 EEGNet 在 8-bit 条件下仅损失约 $0.4\%$ 精度，却获得了最高 $85\%$ 的内存降低和显著的能效提升 \cite{Sch20}。En\'eriz 等在更低成本 MCU 上报告的 PTQ 结果表明，即便在不重新训练量化参数的条件下，INT8 推理仍能将精度损失控制在较小范围内 \cite{Ene23}。Wang 等提出的 MI-BMInet 则将通道选择与 8-bit 量化结合，在保持 $82.51\%$ 二分类准确率的同时，使输入通道数减少约 $6.4$ 倍，并将单次推理能耗压缩到 $30\ \mu\text{J}$ 量级、时延降至 $2.95\ \text{ms}$ \cite{Wan22b}。这些工作说明，对于 MI 脑电这类输入维度有限、窗口长度可控的任务，量化与结构压缩并不是附属优化，而是决定端侧可行性的基础手段。

对于 Transformer 或其他时序模型而言，仅仅降低位宽往往还不够，因为注意力矩阵构造、张量重排和缓存访问模式同样会决定真实运行成本。Busia 等针对 EEGformer 的硬件导向实现表明，当模型规模受控并结合面向 MCU 的实现策略时，Transformer 类结构在 GAP9 上可达到 $13.7\ \text{ms}$ 的单次推理时延和 $0.31\ \text{mJ}$ 的单次能耗 \cite{bus2024reducing}。Jung 等进一步指出，面向低功耗 MCU 的 Tiny Transformer 部署需要同时优化融合权重、深度优先分块和内存布局，其框架相较于常规 ARM / RISC-V 推理库可实现平均 $4.79$ 倍的时延降低与最高 $6.19$ 倍的峰值内存下降 \cite{Jun24}。因此，边缘部署中的模型优化不应被理解为单一的“量化步骤”，而应被视为网络结构、数值表示和运行时算子调度的联合设计问题。

\subsection{平台形态与 RK3568 适配边界}

从平台形态看，当前与脑电端侧部署最相关的硬件大致可分为三类：一是以 Cortex-M、PULP、GAP8 / GAP9 为代表的超低功耗 MCU，优势在于能耗极低、实时性可控，但片上内存与模型规模受限；二是以 Jetson Nano、RK3568 为代表的嵌入式 Linux SBC / NPU 平台，优势在于系统集成度高、接口丰富、适合与采集和通信模块协同工作；三是以中小型 FPGA 为代表的专用加速平台，其优势在于可通过深度定制数据流获得较高的能效比，但开发成本和工具链复杂度通常更高。

不同平台之间的差异并不只体现在峰值算力上，更体现在对同一网络拓扑的执行效率上。Pacini 等对 CPU、Jetson Nano 与 FPGA 上的 MI 分类 CNN 部署进行了横向比较，结果显示 CPU 基线的单次推理时延约为 $121.57\ \text{ms}$、功率约为 $13.22\ \text{W}$、内存占用约为 $476\ \text{MB}$；相较之下，FPGA 方案在保持约 $77\%$ 测试准确率的同时，将单次推理时延降至 $12.97\ \text{ms}$、功率降至 $1.47\ \text{W}$、内存占用降至 $6.37\ \text{MB}$ \cite{Pac24}。这类结果说明，面向边缘部署的模型设计不能简单套用桌面 CPU 上的经验指标，而应结合目标平台的存储层级、整数算子支持和数据搬运代价进行联合评估。

对本文后续采用的 RK3568 平台而言，现有公开证据能够支撑“平台具备部署潜力”这一判断，但尚不足以直接替代本论文的实测验证。You 和 Weng 在 RK3568 SoC 上将时序模型转换为 RKNN 后，报告一段一分钟音频信号的解码时间由 CPU 上的 $12.63\ \text{s}$ 降至 NPU 上的 $1.67\ \text{s}$，表明 RKNN 对时序网络确实能够带来明显加速 \cite{You24}。Ismagilov 的研究则从板级角度给出了 RK3568 在量化 / 半精度模型上的功耗和推理时延量级，其结果显示该类平台具备数百毫秒级推理和数瓦级功耗的实际运行能力 \cite{Ism23}。不过，这两类证据分别来自音频时序任务和视觉目标检测任务，并不能直接等同于“运动想象脑电 Transformer 在 RK3568 上的端到端时延结论”。因此，针对 RK3568 / RKNN 的部署论证应建立在“已有相近平台与相近算子形态的外部证据 + 本文后续实际测试结果”这两个层面之上，而不能仅凭平台名称做出过强推广。

\subsection{蓝牙低功耗（BLE）传输技术}

\subsubsection{BLE 协议栈与低功耗工作机制}

蓝牙低功耗（Bluetooth Low Energy, BLE）工作于 $2.4\ \text{GHz}$ ISM 频段，其设计目标是在牺牲部分持续吞吐能力的前提下，获得更低的平均功耗与更灵活的连接管理机制。BLE 协议栈通常可划分为物理层（PHY）、链路层（Link Layer）、逻辑链路控制与适配层（L2CAP）、属性协议（ATT）和通用属性配置文件（GATT）等部分。对可穿戴采集设备而言，常见工作方式是由设备端作为 Peripheral / GATT Server 暴露数据服务，移动端或边缘主控作为 Central / GATT Client 发起扫描、连接、服务发现与特征值订阅。连接建立后，数据在离散的连接事件（Connection Event）中被分批发送，而非像经典蓝牙那样长时间维持高占空比链路，这也是 BLE 得以降低平均能耗的重要原因。

BLE 的时延与功耗主要受广播间隔、连接间隔、从机延迟和 MTU（Maximum Transmission Unit）等参数影响。连接间隔越短，单位时间内可触发的数据交换次数越多，链路时延更低，但射频唤醒次数上升，平均功耗随之增加；MTU 越大，单个数据包中可承载的有效负载越多，有利于降低协议头开销，但也会增加丢包后的重传代价。在实时信号传输场景中，设备通常将高频连续采样数据放在 Notify / Indicate 类型的特征值中周期性上报，而将采样率、增益、工作模式等控制信息放在可写特征值中完成双向配置。第 3 章中的蓝牙传输电路和固件逻辑，正是围绕这一协议模型展开具体实现。

\subsubsection{BLE 在脑电连续采样中的适配要点}

对脑电采集链路而言，BLE 能否满足实时传输要求，首先取决于原始数据速率与链路有效吞吐之间的匹配关系。若采集通道数为 $C$，每通道采样率为 $f_s$，每个样本在传输层占用 $B_s$ 字节，则原始数据速率近似为
\begin{equation}
R_{\mathrm{data}} = C \cdot f_s \cdot B_s.
\label{eq:ble_rate}
\end{equation}
考虑包头、时间戳、序列号和状态位等附加字段后，链路需要满足 $R_{\mathrm{BLE,eff}} > R_{\mathrm{data}} + R_{\mathrm{meta}}$，其中 $R_{\mathrm{meta}}$ 表示元数据与协议开销对应的等效速率。对双通道或少通道脑电设备而言，只要连接间隔与数据包长度设置合理，BLE 通常能够满足原始波形与少量状态信息的持续传输；但若同时叠加多模态传感器、高频频谱特征或大尺寸控制消息，则链路缓冲和分包策略就会迅速成为系统瓶颈。

除吞吐之外，脑电流式传输还强调时序完整性与低抖动。由于后续解码通常基于固定滑动窗口执行，若包级时间戳不准确、序列号不连续或缓冲区管理不稳定，边缘端接收到的时间序列就会出现错位、突发缺样或窗口边界漂移，进而影响特征提取和模型推理结果。因此，面向脑电的 BLE 传输设计一般需要同时具备四个特性：其一，数据包中显式包含序列号或采样时钟信息，以支持丢包检测与重排序；其二，采用固定长度或近似固定长度的采样批次，便于接收端进行滑动窗口重构；其三，在固件侧保留发送缓存与流控逻辑，以避免采样线程和发包线程相互阻塞；其四，在连接异常或短时链路抖动后支持快速重连与状态恢复。也就是说，BLE 在脑电系统中的角色并不仅是“无线传输通道”，而是连接采样节拍与边缘推理节拍的关键中间层。

\subsection{脑电采集芯片工作原理}

\subsubsection{脑电前端的基本信号链}

脑电信号的幅值通常处于微伏量级，且极易受到工频干扰、电极接触阻抗变化以及人体共模噪声影响，因此采集芯片的首要任务是以尽可能低的附加噪声完成差分放大与模数转换。一个典型的脑电前端信号链通常包括电极接口与保护网络、偏置与参考电路、低噪声差分放大器、可编程增益放大器（PGA）、模拟滤波单元、高分辨率模数转换器（ADC）以及串行数字输出接口等部分。若头皮两点输入电位分别为 $v_+(t)$ 与 $v_-(t)$，系统差分输入可写为
\begin{equation}
v_d(t) = v_+(t) - v_-(t).
\label{eq:eeg_diff_input}
\end{equation}
在增益为 $G$ 的前端放大级之后，进入 ADC 的模拟量可近似表示为
\begin{equation}
v_g(t) = G \cdot v_d(t) + n(t),
\label{eq:eeg_gain_output}
\end{equation}
其中 $n(t)$ 表示由器件热噪声、电源噪声和环境耦合共同构成的等效附加噪声。由式 \eqref{eq:eeg_gain_output} 可见，脑电前端并不是简单地“把信号放大”，而是在尽量抑制共模和外界噪声的前提下，将微弱差分电压映射到 ADC 可分辨的动态范围内。

在离散采样阶段，ADC 以采样周期 $T_s$ 对放大后的信号进行量化。若参考电压为 $V_{\mathrm{ref}}$、ADC 分辨率为 $N$ bit，则其最小量化步长可近似表示为
\begin{equation}
\Delta_{\mathrm{ADC}} = \frac{V_{\mathrm{ref}}}{2^N}.
\label{eq:adc_lsb}
\end{equation}
步长 $\Delta_{\mathrm{ADC}}$ 越小，对微弱脑电波动的分辨能力越强，但同时也要求前级放大和滤波设计更加稳定，否则放大后的干扰同样会被高分辨率采样“完整记录”。因此，脑电采集芯片的性能通常不能脱离输入阻抗、共模抑制能力、增益设置和采样同步性等因素单独评价。

\subsubsection{KS1092 双通道专用前端的工作流程}

结合本论文后续硬件实现，第 3 章所采用的 KS1092 属于面向便携式脑电采集场景的双通道专用模拟前端。其工作流程可以概括为：首先由电极将头皮表面的微弱电位变化转换为差分模拟输入；随后芯片内部或外围的偏置 / 参考网络将输入共模电平稳定在合适范围内，并通过高输入阻抗前端减小电极阻抗波动对信号的直接影响；之后，差分信号经过低噪声放大与可编程增益调节，被映射到 ADC 的有效动态范围内；最后，量化后的数字样本通过串行接口送入主控芯片，由固件完成缓存、时间对齐与无线发送。

对双通道脑电系统而言，采样同步性与通道一致性尤为重要。一方面，运动想象、认知状态识别等任务通常依赖不同通道之间的相对幅值变化、相位关系或窗口内统计特征，若双通道采样时刻不一致，后续时序分析会出现额外误差；另一方面，增益设置、参考稳定性和供电纹波也会影响两个通道之间的可比性。正因如此，便携式脑电专用前端通常强调差分采集、同步采样、低噪声供电和数字接口的稳定输出，而第 3 章中的电路设计、陷波处理与数据缓冲机制，正是围绕上述采集链要求展开的具体工程实现。

\section{本章小结}

本章围绕后续系统实现与算法部署所需的基础技术，对运动想象脑电生理特征、Transformer 编码原理、高效局部注意力机制、边缘部署约束、BLE 无线传输以及脑电采集前端工作流程进行了系统梳理。与绪论中的研究现状不同，本章的重点在于明确后续章节必须依托的客观技术基础：在信号层面，运动想象脑电具有中央区局部化、$\mu/\beta$ 频带分化和显著个体差异；在模型层面，局部窗口注意力能够以更低复杂度建模脑电时序依赖；在系统层面，模型压缩、链路时序和采集前端的一致性共同决定端侧实时性是否成立。

这些基础结论将直接支撑后续章节的展开。第 3 章将在本章给出的 BLE 与采集前端原理基础上完成硬件系统设计；第 4 章则将在本章关于高效注意力和边缘部署约束的分析基础上，进一步构建面向 RK3568 平台的轻量化脑电解码模型。
